%! Project: Design and Conduct a Survey — Unit 1, Lesson 1.8 — Exit Ticket (Answer Key)
\documentclass[10pt]{article}
\usepackage{saar-article}
\usepackage{saar-key}   % pulls in -boxes; do NOT also load -boxes
\newcommand{\TallMath}[1]{$\displaystyle #1\rule[-1.4em]{0pt}{3.2em}$}

\begin{document}

\pageheader{Unit 1, Lesson 1.8}{Exit Ticket --- Answer Key}
% NAMESTRIP: no \namedateperiod here — mirrors the blank. Teacher-only prose goes
% in the lesson plan as \begin{teachernote}[Exit Ticket], never in a key.

\vspace{0.15in}

\begin{notesbox}{Exit Ticket --- No Notes}
\small
A homeroom of $30$ students at \textbf{Bayside Middle School} surveyed itself
about the spirit-week theme and reported the result to the whole school of
$750$. Every student in the homeroom answered.

\vspace{3pt}
\begin{center}
\renewcommand{\arraystretch}{1.3}
\begin{tabularx}{0.92\linewidth}{@{} X c c c c @{}}
\toprule
\textbf{Theme} & Decades & Color wars & Movie characters & Sports teams \\
\midrule
\textbf{Students} & $12$ & $9$ & $6$ & $3$ \\
\bottomrule
\end{tabularx}
\end{center}

\begin{enumerate}[leftmargin=1.6em, itemsep=8pt, topsep=4pt]

  \item Work with the homeroom's data.

        \textbf{(a)} What percent of the $30$ chose decades?

        \begin{work}
          \dfrac{12}{30} &= 0.40 \\
                         &= 40\%
        \end{work}

        \textbf{(b)} If that held school-wide, about how many of the $750$
        students would choose decades?

        \begin{work}
          0.40 \times 750 &= 300 \text{ students}
        \end{work}

  \item The homeroom wants to display its result.

        \textbf{(a)} ``Theme'' is a \ans{categorical} variable, so the display
        that fits is a \ans{bar graph}.

        \textbf{(b)} On that display, what does the height of the tallest bar
        tell a reader?
        \ansline{The share of the homeroom that chose decades --- $40\%$, the most common answer.}

  \item The homeroom's report says: \emph{``Bayside students want a decades
        theme.''}

        \textbf{(a)} Name the flaw in that sentence and rewrite it so it is
        honest.
        \ansline{It claims all $750$ students from a convenience sample of one homeroom. Honest version: ``Of the $30$ students in this homeroom, $40\%$ chose a decades theme.''}

        \textbf{(b)} What one change to the \emph{sample} would have earned the
        original wording?
        \ansline{Choose the students by chance from all $750$ --- for example a stratified sample across the grades.}

\end{enumerate}
\end{notesbox}

\end{document}
